\documentclass[11pt]{article}
\usepackage[T1]{fontenc}
\usepackage{lmodern,microtype,amsmath,amssymb,amsthm,booktabs,longtable,array}
\usepackage[margin=1in]{geometry}
\usepackage{xurl}
\usepackage[hidelinks]{hyperref}
\setlength{\emergencystretch}{3em}
\newtheorem{theorem}{Theorem}
\newtheorem{lemma}[theorem]{Lemma}
\newtheorem{proposition}[theorem]{Proposition}
\providecommand{\tightlist}{\setlength{\itemsep}{0pt}\setlength{\parskip}{0pt}}
\providecommand{\passthrough}[1]{#1}
\hypersetup{pdfauthor={Yongxian Zhang}}

\title{A counterexample to the prime-order-factor conjecture for unimodal cyclotomic generating functions}
\author{Yongxian Zhang\\[2pt]\small School of Computer Science and Engineering\\\small South China University of Technology\\\small\texttt{mxymmxym1@gmail.com}\\\small ORCID: 0009-0000-3864-3536}
\date{8 October 2026}
\begin{document}\maketitle
\begin{abstract}We exhibit an explicit unimodal basic cyclotomic generating function of degree 216 whose cyclotomic factors have only the composite orders 4, 9, 25 and 30. A complete positive centered q-integer expansion certifies strict unimodality. The example refutes the prime-order-factor conjecture of Billey and Swanson. The actual polynomial, coefficient certificate, basic quotient identity and universal prime exclusion have complete Lean proofs.\end{abstract}
\section{The conjecture and the counterexample}
Billey and Swanson~\cite{BS}, Conjecture 48, propose that every nonconstant
unimodal basic cyclotomic generating function has a prime-order cyclotomic
factor. Here $[m]_q=1+q+\cdots+q^{m-1}$ for $m\ge1$; a basic CGF is a
polynomial with nonnegative integer coefficients that is a quotient of
products of these $q$-integers, without an additional scalar or monomial
factor. The conjecture imposes no squarefreeness condition or exclusion
of powers. The predicted factor is $\Phi_p$ for a prime $p$, not an
arbitrary prime-power order.

\begin{theorem}
The polynomial
\[
 F(q)=\bigl(\Phi_4(q)\Phi_9(q)\Phi_{25}(q)\Phi_{30}(q)\bigr)^6
\]
is a nonconstant basic CGF of degree 216. Its integer coefficients are
strictly positive, strictly increasing up to degree 108 and strictly
decreasing thereafter. No $\Phi_p$ with $p$ prime divides it.
Consequently Conjecture 48 of~\cite{BS} is false.
\end{theorem}

\section{Factorization and the basic-CGF identity}
Put $B=\Phi_4\Phi_9\Phi_{25}\Phi_{30}$. The factors are
\begin{align*}
 \Phi_4(q)&=1+q^2,\\
 \Phi_9(q)&=1+q^3+q^6,\\
 \Phi_{25}(q)&=1+q^5+q^{10}+q^{15}+q^{20},\\
 \Phi_{30}(q)&=1+q-q^3-q^4-q^5+q^7+q^8.
\end{align*}
The identity $q^m-1=\prod_{d\mid m}\Phi_d(q)$ implies
\begin{equation}\label{eq:quotient}
 B(q)=\frac{[4]_q[9]_q[25]_q[30]_q}
                    {[1]_q[6]_q[10]_q[15]_q}.
\end{equation}
Indeed, after cancelling the four copies of $q-1$ on each side, the
cyclotomic factors of the numerator are those with orders dividing
4, 9, 25 or 30, with multiplicities. Cancelling the factors whose orders
divide 1, 6, 10 or 15 leaves exactly 4, 9, 25 and 30 once each.
The denominator is a nonzero polynomial. Raising \eqref{eq:quotient} to
the sixth power supplies the required quotient representation for $F$.
The remaining requirement of coefficient nonnegativity is verified next.

\section{A complete finite coefficient certificate}
All four factors are reciprocal. Consequently if
$F=\sum_{j=0}^{216}a_jq^j$, then $a_j=a_{216-j}$.
Exact integer multiplication yields
\begin{equation}\label{eq:centered}
 F(q)=\sum_{j=0}^{108}d_jq^j[217-2j]_q.
\end{equation}
The complete list of weights follows; each row starts at its displayed
index and lists consecutive values. The last row has only four values.
\begin{longtable}{r|rrrrr}
$j$ & \multicolumn{5}{c}{Consecutive values $d_j,d_{j+1},\ldots$}\\\hline
0 & 1 & 5 & 15 & 35 & 64\\
5 & 96 & 126 & 150 & 180 & 238\\
10 & 326 & 426 & 501 & 531 & 567\\
15 & 701 & 991 & 1371 & 1692 & 1866\\
20 & 1980 & 2326 & 3071 & 4023 & 4806\\
25 & 5124 & 5250 & 5856 & 7269 & 9147\\
30 & 10585 & 11189 & 11655 & 13079 & 15982\\
35 & 19290 & 21585 & 22779 & 24183 & 27587\\
40 & 32545 & 36825 & 39095 & 40213 & 43527\\
45 & 50235 & 57570 & 62100 & 62948 & 64882\\
50 & 72555 & 84635 & 95020 & 97758 & 96567\\
55 & 100479 & 113595 & 130975 & 140138 & 138138\\
60 & 135060 & 142464 & 162645 & 180663 & 184173\\
65 & 177591 & 176555 & 194149 & 218922 & 230692\\
70 & 223826 & 212586 & 221560 & 248612 & 269481\\
75 & 265241 & 241396 & 232836 & 254723 & 286675\\
80 & 295956 & 268346 & 238156 & 241074 & 275394\\
85 & 304410 & 288381 & 244691 & 218995 & 235689\\
90 & 273250 & 275390 & 232581 & 184211 & 174619\\
95 & 208125 & 227485 & 197435 & 138246 & 102744\\
100 & 124275 & 158475 & 150360 & 91356 & 28794\\
105 & 23500 & 59600 & 78150 & 38406 &\\
\end{longtable}

Identity~\eqref{eq:centered} is a finite polynomial certificate: expanding
the four displayed factors and raising their product to power six gives
its left-hand side, and expanding the 109 specified centered $q$-integers
gives its right-hand side. The coefficient list and an independent
integer-arithmetic verifier of both sides accompany this paper; the
complete identity is also proved in Lean.
For $0\le k\le108$, a term with index $j$ contributes to degree $k$
exactly when $j\le k$: its support is $[j,216-j]$. Hence
\[
 a_k=\sum_{j=0}^k d_j.
\]
Every listed $d_j$ is strictly positive, so the first-half coefficients
are strictly positive and strictly increasing. Reciprocity gives the
strictly decreasing second half. This proves the required nonnegativity
and strict unimodality, with a unique peak at 108. The coefficient sum
is $30^6=729000000$ and $a_{108}=11434392$.

\section{Excluding every prime-order factor}
Cyclotomic polynomials are irreducible over $\mathbb Q$, and distinct
cyclotomic polynomials are relatively prime. The complete factorization
of $F$ has only orders 4, 9, 25 and 30, none prime. Thus $\Phi_p$ does not
divide $F$ for any prime $p$. Equivalently, a primitive $p$th root cannot
have any of these four distinct composite orders. This proves the
universal prime exclusion without extrapolating from finitely many primes.
Finally $\deg F=6(2+6+20+8)=216>0$, so $F\ne1$. All assertions of the
theorem follow.

\section{Reproduction and scope}
Run \texttt{python3 verify\_counterexample.py} in this paper's directory.
The standard-library checker reconstructs the cyclotomic factors from
$q^m-1$, verifies \eqref{eq:quotient}, multiplies out all 217 coefficients,
checks all 109 weights and independently reconstructs \eqref{eq:centered}.
It also checks exact remainders for every possible prime up to 217; larger
prime-order divisors are excluded by degree. That finite check supplements
the universal factorization proof above. There is no numerical-root
computation. No minimum-degree or worldwide-first claim is made, and the
distinct later HSOP-membership problem is not the conjecture addressed here.
\section{Formal verification and reproduction}
The principal theorem has a complete Lean proof. The precise scope is:
The explicit degree-216 polynomial, strict unimodality, basic q-integer quotient, and exclusion of every prime-order cyclotomic factor. No degree-minimality claim.
The public proof source snapshot is fixed at commit
\nolinkurl{02af7db99db789fecb0d5f5f8f37b1a4e40762a0} in \url{https://github.com/mxym/math}.
The proof packages, with their own dependency locks, audit records and
reproduction instructions, are:
\begin{itemize}
\item \nolinkurl{formalizations/cyclotomic-prime-factor-counterexample}.
\end{itemize}
The main formal entry points are:
\begin{itemize}
\item \nolinkurl{CyclotomicCounterexample.explicit_counterexample}.
\item \nolinkurl{CyclotomicCounterexample.conjecture48_false}.
\item \nolinkurl{CyclotomicCounterexample.F_centered_certificate}.
\item \nolinkurl{CyclotomicCounterexample.F_qInteger_quotient_certificate}.
\end{itemize}
Lean is fixed to version 4.34.1, compiler commit
\nolinkurl{5045d0056413266e57c625dcd7c365b10e377c52}; Mathlib is fixed to
\nolinkurl{d13f23b723b8a846827a245b89c10fc7d3f11612}.
The audited mathematical roots depend only on the standard axioms
\texttt{propext}, \texttt{Classical.choice}, and \texttt{Quot.sound}.
The archived audits state their actual compilation and replay scopes;
publication of this typeset version does not imply a new execution of every
previous empty-kernel replay. The companion README maps each printed
principal result to its formal theorem and explains how to repeat verification.

\section*{Research and preparation disclosures}
This research, formalization and manuscript preparation used AI assistance.
The author is responsible for the manuscript. The principal claims are
supported by the readable proofs and the public Lean artifacts specified
above. Formal verification is not external human peer review; no such review
is claimed. No external funding was received. No worldwide priority or
journal acceptance is asserted.
\begin{thebibliography}{9}
\bibitem{BS} S. C. Billey and J. P. Swanson,
\emph{Cyclotomic Generating Functions}, Electronic Journal of Combinatorics
\textbf{31}(4) (2024), P4.4. \url{https://doi.org/10.37236/12687}.
Pinned manuscript: \url{https://arxiv.org/abs/2305.07620v3},
Definitions 1, 3 and 45; Conjecture 48, printed page 32.
\end{thebibliography}
\end{document}
